% ============================================================
\chapter{Fixed Points in Ordered Spaces}
\label{ch:ordered}
% ============================================================

In 1928, Bronislaw Knaster and Alfred Tarski proved a result of extraordinary simplicity and reach: every monotone map on a complete lattice has a fixed point. Better still, the set of \emph{all} fixed points itself forms a complete lattice. This theorem uses no metric, no topology, no continuity: only the order structure intervenes. It would find spectacular applications in theoretical computer science, where it underpins the denotational semantics of recursive programs (Dana Scott's work in the 1970s), and in logic, where it supports the theory of inductive definitions. Together with the Bourbaki--Witt transfinite iterates and Kleene's theorem for $\omega$-cpo's, this chapter explores the third great family of fixed point theorems --- the one founded on order.

\begin{intuition}
Fixed point theorems in ordered sets constitute the third great family of results, after
the metric theorems (Banach) and the topological ones (Brouwer, Schauder). Here, no metric
or topological structure is required: only the order structure intervenes. The fundamental
result is the Knaster--Tarski theorem, which asserts that every monotone map on a complete
lattice has a fixed point, and that the set of fixed points itself forms a complete lattice.
\end{intuition}

% ------------------------------------------------------------
\section{Background on Ordered Sets}
% ------------------------------------------------------------

\begin{definition}[Partial order and lattice]
\index{partial order}
A \emph{partially ordered set} (poset) is a pair $(P, \leq)$ where $\leq$ is a reflexive,
antisymmetric, and transitive relation. A \emph{lattice} is a poset in which every pair
$\{a, b\}$ has a supremum $a \vee b$ and an infimum $a \wedge b$. A lattice is
\emph{complete} if every subset has a supremum and an infimum.
\end{definition}

\begin{example}[Fundamental complete lattices]
\begin{itemize}
  \item $(\mathcal{P}(X), \subseteq)$ for any set $X$.
  \item $([0,1], \leq)$ with the usual order.
  \item The lattice of vector subspaces of a vector space.
  \item The lattice of equivalence relations on a set.
\end{itemize}
\end{example}

\begin{definition}[Monotone and inflationary maps]
\index{monotone map}
Let $(P, \leq)$ be a poset. A map $f\colon P \to P$ is \emph{monotone} (or
order-preserving) if $x \leq y \Rightarrow f(x) \leq f(y)$. It is \emph{inflationary} if
$x \leq f(x)$ for all $x$, and \emph{deflationary} if $f(x) \leq x$ for all $x$.
\end{definition}

% ------------------------------------------------------------
\section{The Knaster--Tarski Theorem}
% ------------------------------------------------------------

\begin{theorem}[Knaster--Tarski]
\label{thm:knaster-tarski}
\index{Knaster--Tarski theorem}
Let $(L, \leq)$ be a complete lattice and $f\colon L \to L$ a monotone map. Then:
\begin{enumerate}
  \item $f$ has a least fixed point:
    $\mathrm{lfp}(f) = \bigwedge \{x \in L : f(x) \leq x\}$.
  \item $f$ has a greatest fixed point:
    $\mathrm{gfp}(f) = \bigvee \{x \in L : x \leq f(x)\}$.
  \item The set $\mathrm{Fix}(f) = \{x \in L : f(x) = x\}$ is a complete lattice (under
    the induced order).
\end{enumerate}
\end{theorem}

\begin{proof}
Let $\mathrm{Pre}(f) = \{x \in L : f(x) \leq x\}$ (the pre-fixed points) and
$a = \bigwedge \mathrm{Pre}(f)$.

\emph{Step 1:} $f(a) \leq a$. Indeed, for every $x \in \mathrm{Pre}(f)$, we have
$a \leq x$, so $f(a) \leq f(x) \leq x$. Thus $f(a)$ is a lower bound of
$\mathrm{Pre}(f)$, hence $f(a) \leq a$.

\emph{Step 2:} $a \leq f(a)$. From Step~1, $f(a) \leq a$, so by monotonicity
$f(f(a)) \leq f(a)$. Thus $f(a) \in \mathrm{Pre}(f)$, hence $a \leq f(a)$.

\emph{Step 3:} $f(a) = a$ by Steps 1 and 2.

\emph{Least fixed point:} if $b$ is a fixed point, then $b \in \mathrm{Pre}(f)$, so
$a \leq b$.

\emph{Complete lattice of fixed points:} Let $S \subseteq \mathrm{Fix}(f)$. Define
$L_S = \{x \in L : x \geq s \; \forall s \in S \text{ and } f(x) \leq x\}$. Then $L_S$
is nonempty ($\top \in L_S$) and $\bigwedge L_S$ is a fixed point of $f$ that is the
supremum of $S$ in $\mathrm{Fix}(f)$. Dually, $\mathrm{Fix}(f)$ admits infima.
\end{proof}

\begin{example}[Application in logic]
Let $X$ be a set and $\Phi\colon \mathcal{P}(X) \to \mathcal{P}(X)$ a monotone operator.
The least fixed point of $\Phi$ is the smallest set $A$ with $\Phi(A) = A$. If $\Phi$ is
a logical consequence operator, $\mathrm{lfp}(\Phi)$ is the set of deducible formulas.
\end{example}

% ------------------------------------------------------------
\section{The Bourbaki--Witt Theorem}
% ------------------------------------------------------------

\begin{theorem}[Bourbaki--Witt]
\label{thm:bourbaki-witt}
\index{Bourbaki--Witt theorem}
Let $(P, \leq)$ be a partially ordered set in which every chain has an upper bound, and
let $f\colon P \to P$ be an inflationary map ($x \leq f(x)$ for all $x$). If $P$ has a
least element $\bot$, then $f$ has a fixed point.
\end{theorem}

\begin{proof}
Consider the collection $\mathcal{A}$ of all chains $C$ in $P$ such that:
\begin{itemize}
  \item $\bot \in C$;
  \item $C$ is closed under $f$;
  \item if $C' \subseteq C$ is a chain, any upper bound of $C'$ in $P$ that belongs to
    $C$ is indeed an upper bound.
\end{itemize}
Define $C_0 = \bigcap_{C \in \mathcal{A}} C$. One shows that $C_0$ is a well-ordered
chain. Let $m$ be an upper bound of $C_0$ in $P$ (which exists by hypothesis). Then
$f(m) \geq m$. If $f(m) = m$, it is a fixed point. Otherwise, $C_0 \cup \{m\}$ would be a
strictly larger chain, contradicting the minimality of $C_0$.
\end{proof}

\begin{remark}
The Bourbaki--Witt theorem is equivalent to Zorn's lemma in ZF + axiom of choice. Without
the axiom of choice, the statement remains true for explicitly defined inflationary
functions.
\end{remark}

% ------------------------------------------------------------
\section{Application to Formal Concept Analysis}
% ------------------------------------------------------------

\begin{definition}[Formal context]
\index{formal concept analysis}
A \emph{formal context} is a triple $(G, M, I)$ where $G$ is a set of objects, $M$ a set
of attributes, and $I \subseteq G \times M$ an incidence relation ($gIm$ means that
object $g$ has attribute $m$).
\end{definition}

\begin{definition}[Derivation operators]
For $A \subseteq G$ and $B \subseteq M$, define:
\[
  A' = \{m \in M : \forall g \in A, \; gIm\}, \qquad
  B' = \{g \in G : \forall m \in B, \; gIm\}.
\]
\end{definition}

\begin{proposition}[Galois connection]
The operators $A \mapsto A'$ and $B \mapsto B'$ form a \emph{Galois connection} between
$(\mathcal{P}(G), \subseteq)$ and $(\mathcal{P}(M), \supseteq)$:
\begin{enumerate}
  \item $A_1 \subseteq A_2 \Rightarrow A_2' \subseteq A_1'$;
  \item $A \subseteq A''$;
  \item $A' = A'''$.
\end{enumerate}
\end{proposition}

\begin{definition}[Formal concept]
A \emph{formal concept} is a pair $(A, B)$ with $A \subseteq G$, $B \subseteq M$,
$A' = B$ and $B' = A$. The set $A$ is the \emph{extent} and $B$ the \emph{intent} of the
concept.
\end{definition}

\begin{theorem}[Concept lattice]
\label{thm:concept-lattice}
The set of formal concepts of a context $(G, M, I)$, ordered by
$(A_1, B_1) \leq (A_2, B_2) \Leftrightarrow A_1 \subseteq A_2$ (equivalently
$B_2 \subseteq B_1$), is a complete lattice.
\end{theorem}

\begin{proof}
The closure operator $\varphi\colon \mathcal{P}(G) \to \mathcal{P}(G)$ defined by
$\varphi(A) = A''$ is monotone, extensive ($A \subseteq A''$), and idempotent
($A'' = A''''$). The closed sets (fixed points of $\varphi$) form a complete lattice and
correspond bijectively to formal concepts.
\end{proof}

% ------------------------------------------------------------
\section{Fixed Points and Closure Operators}
% ------------------------------------------------------------

\begin{definition}[Closure operator]
\index{closure operator}
A \emph{closure operator} on a poset $(P, \leq)$ is a monotone, extensive, and idempotent
map $\varphi\colon P \to P$.
\end{definition}

\begin{theorem}
Let $(L, \leq)$ be a complete lattice and $\varphi\colon L \to L$ a closure operator. Then
$\mathrm{Fix}(\varphi)$ is a complete lattice (with suprema of $L$ closed by $\varphi$,
and infima unchanged).
\end{theorem}

\begin{proof}
The fixed points of $\varphi$ are the closed elements, i.e., those $x$ with
$\varphi(x) = x$. For $S \subseteq \mathrm{Fix}(\varphi)$, the infimum in
$\mathrm{Fix}(\varphi)$ is $\bigwedge_L S$ (since $\varphi(\bigwedge_L S) \leq
\varphi(s) = s$ for all $s \in S$, hence $\varphi(\bigwedge_L S) \leq \bigwedge_L S$,
and by extensiveness we get equality). The supremum in $\mathrm{Fix}(\varphi)$ is
$\varphi(\bigvee_L S)$.
\end{proof}

\begin{keyformulas}[title=\textbf{Fixed points in ordered sets}]
\begin{center}
\renewcommand{\arraystretch}{1.3}
\begin{tabular}{|l|l|l|}
\hline
\textbf{Theorem} & \textbf{Hypotheses} & \textbf{Conclusion} \\
\hline
Knaster--Tarski & Complete lattice, $f$ monotone & $\mathrm{Fix}(f)$ complete lattice \\
Bourbaki--Witt & Chains bounded, $f$ inflat. & Fixed point exists \\
Closure op. & Complete lattice, $\varphi$ closure & $\mathrm{Fix}(\varphi)$ complete lattice \\
\hline
\end{tabular}
\end{center}
\end{keyformulas}

% ------------------------------------------------------------
\section{Exercises}
% ------------------------------------------------------------

\begin{exercise}
Show that the Knaster--Tarski theorem implies the Cantor--Bernstein theorem: if
$f\colon A \hookrightarrow B$ and $g\colon B \hookrightarrow A$ are injections, then
there exists a bijection $h\colon A \to B$.
\emph{Hint: apply the theorem to $\varphi(S) = A \setminus g(B \setminus f(S))$ on
$(\mathcal{P}(A), \subseteq)$.}
\end{exercise}

\begin{exercise}
Let $(L, \leq)$ be a complete lattice and $f\colon L \to L$ monotone. Show that if
$a \leq f(a)$ and $f(b) \leq b$ with $a \leq b$, then $f$ has a fixed point in the
interval $[a, b]$.
\end{exercise}

\begin{exercise}
Give an example of an ordered set (not complete) and a monotone map with no fixed point.
\end{exercise}

\begin{exercise}
Let $G = (\{1,2,3,4\}, E)$ be a graph and
$f\colon \mathcal{P}(\{1,2,3,4\}) \to \mathcal{P}(\{1,2,3,4\})$ defined by
$f(S) = S \cup \{v : \exists u \in S, \; (u,v) \in E\}$. Compute $\mathrm{lfp}(f)$ and
$\mathrm{gfp}(f)$ for the graph $E = \{(1,2), (2,3), (3,4)\}$ starting from $\{1\}$.
\end{exercise}

\begin{exercise}[Galois connections and fixed points]
\index{Galois connection}
Let $(f, g)$ be a Galois connection between two complete lattices $L_1$ and $L_2$
($f \dashv g$). Show that $g \circ f$ is a closure operator on $L_1$ and $f \circ g$ is a
kernel operator on $L_2$.
\end{exercise}

\begin{exercise}
Construct the concept lattice of the formal context with
$G = \{\text{dog}, \text{cat}, \text{fish}\}$,
$M = \{\text{fur}, \text{swims}, \text{4 legs}\}$ and the natural incidence relation.
\end{exercise}
